\section{Part Twelve: The Topology of Error Propagation}

The five functions from Part Four describe what happens \emph{at} a step. They say very little about what happens \emph{between} steps --- and error propagation is fundamentally a between-steps phenomenon. This part treats a workflow as an explicit \textbf{graph of state transformations}, not a pipeline, because a pipeline model quietly assumes information moves cleanly from one stage to the next, and that assumption is exactly where propagation failures hide.

\textbf{Each edge in that graph carries more than content.} Alongside whatever substantive material is being passed along, an edge can carry --- or fail to carry --- evidence, provenance, permissions, uncertainty, assumptions, authority, timestamps, and scope conditions. A genuinely distinct and easy-to-miss form of propagation failure follows directly from this: \textbf{a workflow edge may transmit content correctly while corrupting its provenance, uncertainty, permissions, or scope.} The words that reach the next step can be exactly right while everything \emph{about} those words that a downstream consumer would need to evaluate them properly --- where they came from, how confident anyone should be, who's allowed to see them, what they don't cover --- is silently lost or altered.

\textbf{At any node or handoff, an error can be one of at least eight things}, and these are worth distinguishing because each implies a different fix: \textbf{introduced} (new at this step); \textbf{preserved} unchanged from upstream; \textbf{amplified} (one bad extraction propagating into many downstream records); \textbf{attenuated} (a human or a later step partially correcting it); \textbf{masked} (a polished summary concealing the uncertainty that was actually present in the underlying material); \textbf{converted} (probabilistic advice becoming a binary workflow flag, discarding the very uncertainty a careful reader would have wanted); \textbf{compensated for accidentally} (a second, independent error happening to cancel the first, which is good luck, not a control); and \textbf{embedded in state} (persisting forward to influence future runs, not just the current one).

\textbf{Four propagation shapes are worth naming explicitly, because they call for different mitigations.} \emph{Serial propagation}: an upstream error becomes an unexamined premise for every step downstream of it --- this is the chain-propagation problem Part Nine already raised, now given its general name. \emph{Fan-out amplification}: a single mistaken intermediate result gets copied into many outputs, decisions, or records at once, so the blast radius of one bad step is multiplied by however many places it gets reused. \emph{Feedback propagation}: an output re-enters the system later as future context, training data, organizational memory, or precedent --- the system doesn't just make the same mistake twice, it can start treating its own earlier mistake as evidence. \emph{Common-cause propagation}: apparently independent steps fail together because they secretly share the same model family, source corpus, prompt assumptions, retrieval index, or evaluator --- which means several checks that look independent may not actually be testing anything the others didn't already test.

\textbf{Common-cause propagation deserves a specific, verifiable piece of evidence rather than just an assertion, because it directly bears on a point raised more tentatively in Part Four.} A large-scale empirical study (Kim, Garg, Peng, and Garg, accepted to ICML 2025) evaluated error correlation across more than 350 language models on two popular leaderboards and a resume-screening task, and found substantial correlation in model errors --- on one leaderboard dataset, models agreed roughly 60\% of the time specifically in cases where \emph{both} models were wrong. Shared architecture and shared providers were identified as contributing factors, which is unsurprising --- but the more important finding is that \textbf{larger, more accurate models showed highly correlated errors even across genuinely distinct architectures and different providers}. This is a meaningfully stronger and more specific claim than the general "a checker built from a similar model may share the writer's blind spots" caution raised in Part Four, which was about a single model checking its own or a similar model's work. This finding is about \emph{independently built, differently architected} models from \emph{different labs} still agreeing on the same wrong answers, more than chance would predict, and increasingly so as capability rises. The practical implication: \textbf{using two or more different models to cross-check an answer is real evidence, but weaker evidence than it intuitively feels like}, especially for exactly the frontier-capability cases where a team might be most tempted to treat multi-model agreement as strong reassurance. This finding was measured on leaderboard benchmarks and a hiring task specifically --- extending its magnitude to other domains, including the professional-judgment tasks this framework is chiefly concerned with, is a reasonable inference, not something the cited study establishes directly, and is flagged as such in the debate section at the end of this document.

\textbf{A worked example, to make the abstract propagation shapes concrete.} Consider an AI-assisted contract review. An OCR step silently drops a page. Retrieval, working from what it was actually given, never sees the change-of-control clause that page contained. The model, working correctly from the context it was actually handed, concludes that no such restriction exists --- a conclusion that is entirely reasonable given what it saw. A citation checker verifies that every citation the model \emph{did} make is well-supported, and passes the output, because it was never built to notice a clause that never appeared at all. The interface displays a clean "no issue" flag. A human reviewer, seeing no cited exception and a clean flag, accepts the conclusion --- reasonably, since nothing in front of them suggested a problem. The result is entered into the matter summary, and later into the firm's precedent collection, where it will inform how similar-looking future matters get handled. Walking this through the vocabulary above: the original failure is an \textbf{omission}, not an error in anything that was actually said. It is \textbf{masked} by a conclusion that looks, and in a narrow sense is, entirely valid. The verification step has \textbf{zero coverage} of the actual failure --- it was checking claim validity, not set completeness (a distinction developed further in Part Thirteen). The binary interface \textbf{erases} the underlying uncertainty about page count and document completeness. The human reviewer's acceptance introduces a \textbf{reliance pathway} that a differently designed interface might have interrupted. The error then \textbf{persists in organizational state}, and future retrieval against the precedent collection creates a \textbf{feedback loop} --- later matters may now be evaluated partly against this one's mistaken conclusion. Under an output-only analysis of the kind Part Nine describes, this system looks grounded and internally coherent at every step. Under the fuller topology described here, it's a straightforward, traceable propagation failure from one dropped page to a piece of organizational precedent --- and the eventual harm, whenever it surfaces, may be far larger than the original extraction error would suggest on its own.
